\documentclass[12pt]{article}
\usepackage[canadien]{babel} 
\usepackage[utf8]{inputenc}
\usepackage[T1]{fontenc}
\usepackage{fancyhdr}
\usepackage{graphicx}
\usepackage{enumerate}
\usepackage{amsmath}
\usepackage{amssymb}
\usepackage{subfigure}
\usepackage{amsmath}
\usepackage{amssymb}
\usepackage{bm}
\usepackage{array}
\usepackage{url}
\usepackage[dvipsnames]{xcolor}
\usepackage{todonotes}

\setlength{\parindent}{0cm}
\addtolength{\oddsidemargin}{-2cm}
\addtolength{\evensidemargin}{-2cm}
\setlength{\textwidth}{17.78cm}
\addtolength{\topmargin}{-2.25cm}
\setlength{\textheight}{23.24cm}
\addtolength{\parskip}{5mm}
\pagestyle{fancy}

%************
%* COMMANDS *
%************

%%%%% NEW MATH DEFINITIONS %%%%%

% Mark sections of captions for referring to divisions of figures
\newcommand{\figleft}{{\em (Left)}}
\newcommand{\figcenter}{{\em (Center)}}
\newcommand{\figright}{{\em (Right)}}
\newcommand{\figtop}{{\em (Top)}}
\newcommand{\figbottom}{{\em (Bottom)}}
\newcommand{\captiona}{{\em (a)}}
\newcommand{\captionb}{{\em (b)}}
\newcommand{\captionc}{{\em (c)}}
\newcommand{\captiond}{{\em (d)}}

% Highlight a newly defined term
\newcommand{\newterm}[1]{{\bf #1}}


% Figure reference, lower-case.
\def\figref#1{figure~\ref{#1}}
% Figure reference, capital. For start of sentence
\def\Figref#1{Figure~\ref{#1}}
\def\twofigref#1#2{figures \ref{#1} and \ref{#2}}
\def\quadfigref#1#2#3#4{figures \ref{#1}, \ref{#2}, \ref{#3} and \ref{#4}}
% Section reference, lower-case.
\def\secref#1{section~\ref{#1}}
% Section reference, capital.
\def\Secref#1{Section~\ref{#1}}
% Reference to two sections.
\def\twosecrefs#1#2{sections \ref{#1} and \ref{#2}}
% Reference to three sections.
\def\secrefs#1#2#3{sections \ref{#1}, \ref{#2} and \ref{#3}}
% Reference to an equation, lower-case.
\def\eqref#1{equation~\ref{#1}}
% Reference to an equation, upper case
\def\Eqref#1{Equation~\ref{#1}}
% A raw reference to an equation---avoid using if possible
\def\plaineqref#1{\ref{#1}}
% Reference to a chapter, lower-case.
\def\chapref#1{chapter~\ref{#1}}
% Reference to an equation, upper case.
\def\Chapref#1{Chapter~\ref{#1}}
% Reference to a range of chapters
\def\rangechapref#1#2{chapters\ref{#1}--\ref{#2}}
% Reference to an algorithm, lower-case.
\def\algref#1{algorithm~\ref{#1}}
% Reference to an algorithm, upper case.
\def\Algref#1{Algorithm~\ref{#1}}
\def\twoalgref#1#2{algorithms \ref{#1} and \ref{#2}}
\def\Twoalgref#1#2{Algorithms \ref{#1} and \ref{#2}}
% Reference to a part, lower case
\def\partref#1{part~\ref{#1}}
% Reference to a part, upper case
\def\Partref#1{Part~\ref{#1}}
\def\twopartref#1#2{parts \ref{#1} and \ref{#2}}

\def\ceil#1{\lceil #1 \rceil}
\def\floor#1{\lfloor #1 \rfloor}
\def\1{\bm{1}}
\newcommand{\train}{\mathcal{D}}
\newcommand{\valid}{\mathcal{D_{\mathrm{valid}}}}
\newcommand{\test}{\mathcal{D_{\mathrm{test}}}}

\def\eps{{\epsilon}}


% Random variables
\def\reta{{\textnormal{$\eta$}}}
\def\ra{{\textnormal{a}}}
\def\rb{{\textnormal{b}}}
\def\rc{{\textnormal{c}}}
\def\rd{{\textnormal{d}}}
\def\re{{\textnormal{e}}}
\def\rf{{\textnormal{f}}}
\def\rg{{\textnormal{g}}}
\def\rh{{\textnormal{h}}}
\def\ri{{\textnormal{i}}}
\def\rj{{\textnormal{j}}}
\def\rk{{\textnormal{k}}}
\def\rl{{\textnormal{l}}}
% rm is already a command, just don't name any random variables m
\def\rn{{\textnormal{n}}}
\def\ro{{\textnormal{o}}}
\def\rp{{\textnormal{p}}}
\def\rq{{\textnormal{q}}}
\def\rr{{\textnormal{r}}}
\def\rs{{\textnormal{s}}}
\def\rt{{\textnormal{t}}}
\def\ru{{\textnormal{u}}}
\def\rv{{\textnormal{v}}}
\def\rw{{\textnormal{w}}}
\def\rx{{\textnormal{x}}}
\def\ry{{\textnormal{y}}}
\def\rz{{\textnormal{z}}}

% Random vectors
\def\rvepsilon{{\mathbf{\epsilon}}}
\def\rvtheta{{\mathbf{\theta}}}
\def\rva{{\mathbf{a}}}
\def\rvb{{\mathbf{b}}}
\def\rvc{{\mathbf{c}}}
\def\rvd{{\mathbf{d}}}
\def\rve{{\mathbf{e}}}
\def\rvf{{\mathbf{f}}}
\def\rvg{{\mathbf{g}}}
\def\rvh{{\mathbf{h}}}
\def\rvu{{\mathbf{i}}}
\def\rvj{{\mathbf{j}}}
\def\rvk{{\mathbf{k}}}
\def\rvl{{\mathbf{l}}}
\def\rvm{{\mathbf{m}}}
\def\rvn{{\mathbf{n}}}
\def\rvo{{\mathbf{o}}}
\def\rvp{{\mathbf{p}}}
\def\rvq{{\mathbf{q}}}
\def\rvr{{\mathbf{r}}}
\def\rvs{{\mathbf{s}}}
\def\rvt{{\mathbf{t}}}
\def\rvu{{\mathbf{u}}}
\def\rvv{{\mathbf{v}}}
\def\rvw{{\mathbf{w}}}
\def\rvx{{\mathbf{x}}}
\def\rvy{{\mathbf{y}}}
\def\rvz{{\mathbf{z}}}

% Elements of random vectors
\def\erva{{\textnormal{a}}}
\def\ervb{{\textnormal{b}}}
\def\ervc{{\textnormal{c}}}
\def\ervd{{\textnormal{d}}}
\def\erve{{\textnormal{e}}}
\def\ervf{{\textnormal{f}}}
\def\ervg{{\textnormal{g}}}
\def\ervh{{\textnormal{h}}}
\def\ervi{{\textnormal{i}}}
\def\ervj{{\textnormal{j}}}
\def\ervk{{\textnormal{k}}}
\def\ervl{{\textnormal{l}}}
\def\ervm{{\textnormal{m}}}
\def\ervn{{\textnormal{n}}}
\def\ervo{{\textnormal{o}}}
\def\ervp{{\textnormal{p}}}
\def\ervq{{\textnormal{q}}}
\def\ervr{{\textnormal{r}}}
\def\ervs{{\textnormal{s}}}
\def\ervt{{\textnormal{t}}}
\def\ervu{{\textnormal{u}}}
\def\ervv{{\textnormal{v}}}
\def\ervw{{\textnormal{w}}}
\def\ervx{{\textnormal{x}}}
\def\ervy{{\textnormal{y}}}
\def\ervz{{\textnormal{z}}}

% Random matrices
\def\rmA{{\mathbf{A}}}
\def\rmB{{\mathbf{B}}}
\def\rmC{{\mathbf{C}}}
\def\rmD{{\mathbf{D}}}
\def\rmE{{\mathbf{E}}}
\def\rmF{{\mathbf{F}}}
\def\rmG{{\mathbf{G}}}
\def\rmH{{\mathbf{H}}}
\def\rmI{{\mathbf{I}}}
\def\rmJ{{\mathbf{J}}}
\def\rmK{{\mathbf{K}}}
\def\rmL{{\mathbf{L}}}
\def\rmM{{\mathbf{M}}}
\def\rmN{{\mathbf{N}}}
\def\rmO{{\mathbf{O}}}
\def\rmP{{\mathbf{P}}}
\def\rmQ{{\mathbf{Q}}}
\def\rmR{{\mathbf{R}}}
\def\rmS{{\mathbf{S}}}
\def\rmT{{\mathbf{T}}}
\def\rmU{{\mathbf{U}}}
\def\rmV{{\mathbf{V}}}
\def\rmW{{\mathbf{W}}}
\def\rmX{{\mathbf{X}}}
\def\rmY{{\mathbf{Y}}}
\def\rmZ{{\mathbf{Z}}}

% Elements of random matrices
\def\ermA{{\textnormal{A}}}
\def\ermB{{\textnormal{B}}}
\def\ermC{{\textnormal{C}}}
\def\ermD{{\textnormal{D}}}
\def\ermE{{\textnormal{E}}}
\def\ermF{{\textnormal{F}}}
\def\ermG{{\textnormal{G}}}
\def\ermH{{\textnormal{H}}}
\def\ermI{{\textnormal{I}}}
\def\ermJ{{\textnormal{J}}}
\def\ermK{{\textnormal{K}}}
\def\ermL{{\textnormal{L}}}
\def\ermM{{\textnormal{M}}}
\def\ermN{{\textnormal{N}}}
\def\ermO{{\textnormal{O}}}
\def\ermP{{\textnormal{P}}}
\def\ermQ{{\textnormal{Q}}}
\def\ermR{{\textnormal{R}}}
\def\ermS{{\textnormal{S}}}
\def\ermT{{\textnormal{T}}}
\def\ermU{{\textnormal{U}}}
\def\ermV{{\textnormal{V}}}
\def\ermW{{\textnormal{W}}}
\def\ermX{{\textnormal{X}}}
\def\ermY{{\textnormal{Y}}}
\def\ermZ{{\textnormal{Z}}}

% Vectors
\def\vzero{{\bm{0}}}
\def\vone{{\bm{1}}}
\def\vmu{{\bm{\mu}}}
\def\vtheta{{\bm{\theta}}}
\def\va{{\bm{a}}}
\def\vb{{\bm{b}}}
\def\vc{{\bm{c}}}
\def\vd{{\bm{d}}}
\def\ve{{\bm{e}}}
\def\vf{{\bm{f}}}
\def\vg{{\bm{g}}}
\def\vh{{\bm{h}}}
\def\vi{{\bm{i}}}
\def\vj{{\bm{j}}}
\def\vk{{\bm{k}}}
\def\vl{{\bm{l}}}
\def\vm{{\bm{m}}}
\def\vn{{\bm{n}}}
\def\vo{{\bm{o}}}
\def\vp{{\bm{p}}}
\def\vq{{\bm{q}}}
\def\vr{{\bm{r}}}
\def\vs{{\bm{s}}}
\def\vt{{\bm{t}}}
\def\vu{{\bm{u}}}
\def\vv{{\bm{v}}}
\def\vw{{\bm{w}}}
\def\vx{{\bm{x}}}
\def\vy{{\bm{y}}}
\def\vz{{\bm{z}}}

% Elements of vectors
\def\evalpha{{\alpha}}
\def\evbeta{{\beta}}
\def\evepsilon{{\epsilon}}
\def\evlambda{{\lambda}}
\def\evomega{{\omega}}
\def\evmu{{\mu}}
\def\evpsi{{\psi}}
\def\evsigma{{\sigma}}
\def\evtheta{{\theta}}
\def\eva{{a}}
\def\evb{{b}}
\def\evc{{c}}
\def\evd{{d}}
\def\eve{{e}}
\def\evf{{f}}
\def\evg{{g}}
\def\evh{{h}}
\def\evi{{i}}
\def\evj{{j}}
\def\evk{{k}}
\def\evl{{l}}
\def\evm{{m}}
\def\evn{{n}}
\def\evo{{o}}
\def\evp{{p}}
\def\evq{{q}}
\def\evr{{r}}
\def\evs{{s}}
\def\evt{{t}}
\def\evu{{u}}
\def\evv{{v}}
\def\evw{{w}}
\def\evx{{x}}
\def\evy{{y}}
\def\evz{{z}}

% Matrix
\def\mA{{\bm{A}}}
\def\mB{{\bm{B}}}
\def\mC{{\bm{C}}}
\def\mD{{\bm{D}}}
\def\mE{{\bm{E}}}
\def\mF{{\bm{F}}}
\def\mG{{\bm{G}}}
\def\mH{{\bm{\text{\textbf{H}}}}}
\def\mI{{\bm{I}}}
\def\mJ{{\bm{J}}}
\def\mK{{\bm{K}}}
\def\mL{{\bm{L}}}
\def\mM{{\bm{M}}}
\def\mN{{\bm{N}}}
\def\mO{{\bm{O}}}
\def\mP{{\bm{P}}}
\def\mQ{{\bm{Q}}}
\def\mR{{\bm{R}}}
\def\mS{{\bm{S}}}
\def\mT{{\bm{T}}}
\def\mU{{\bm{U}}}
\def\mV{{\bm{V}}}
\def\mW{{\bm{W}}}
\def\mX{{\bm{X}}}
\def\mY{{\bm{Y}}}
\def\mZ{{\bm{Z}}}
\def\mBeta{{\bm{\beta}}}
\def\mPhi{{\bm{\Phi}}}
\def\mLambda{{\bm{\Lambda}}}
\def\mSigma{{\bm{\Sigma}}}

% Tensor
\DeclareMathAlphabet{\mathsfit}{\encodingdefault}{\sfdefault}{m}{sl}
\SetMathAlphabet{\mathsfit}{bold}{\encodingdefault}{\sfdefault}{bx}{n}
\newcommand{\tens}[1]{\bm{\mathsfit{#1}}}
\def\tA{{\tens{A}}}
\def\tB{{\tens{B}}}
\def\tC{{\tens{C}}}
\def\tD{{\tens{D}}}
\def\tE{{\tens{E}}}
\def\tF{{\tens{F}}}
\def\tG{{\tens{G}}}
\def\tH{{\tens{H}}}
\def\tI{{\tens{I}}}
\def\tJ{{\tens{J}}}
\def\tK{{\tens{K}}}
\def\tL{{\tens{L}}}
\def\tM{{\tens{M}}}
\def\tN{{\tens{N}}}
\def\tO{{\tens{O}}}
\def\tP{{\tens{P}}}
\def\tQ{{\tens{Q}}}
\def\tR{{\tens{R}}}
\def\tS{{\tens{S}}}
\def\tT{{\tens{T}}}
\def\tU{{\tens{U}}}
\def\tV{{\tens{V}}}
\def\tW{{\tens{W}}}
\def\tX{{\tens{X}}}
\def\tY{{\tens{Y}}}
\def\tZ{{\tens{Z}}}


% Graph
\def\gA{{\mathcal{A}}}
\def\gB{{\mathcal{B}}}
\def\gC{{\mathcal{C}}}
\def\gD{{\mathcal{D}}}
\def\gE{{\mathcal{E}}}
\def\gF{{\mathcal{F}}}
\def\gG{{\mathcal{G}}}
\def\gH{{\mathcal{H}}}
\def\gI{{\mathcal{I}}}
\def\gJ{{\mathcal{J}}}
\def\gK{{\mathcal{K}}}
\def\gL{{\mathcal{L}}}
\def\gM{{\mathcal{M}}}
\def\gN{{\mathcal{N}}}
\def\gO{{\mathcal{O}}}
\def\gP{{\mathcal{P}}}
\def\gQ{{\mathcal{Q}}}
\def\gR{{\mathcal{R}}}
\def\gS{{\mathcal{S}}}
\def\gT{{\mathcal{T}}}
\def\gU{{\mathcal{U}}}
\def\gV{{\mathcal{V}}}
\def\gW{{\mathcal{W}}}
\def\gX{{\mathcal{X}}}
\def\gY{{\mathcal{Y}}}
\def\gZ{{\mathcal{Z}}}

% Sets
\def\sA{{\mathbb{A}}}
\def\sB{{\mathbb{B}}}
\def\sC{{\mathbb{C}}}
\def\sD{{\mathbb{D}}}
% Don't use a set called E, because this would be the same as our symbol
% for expectation.
\def\sF{{\mathbb{F}}}
\def\sG{{\mathbb{G}}}
\def\sH{{\mathbb{H}}}
\def\sI{{\mathbb{I}}}
\def\sJ{{\mathbb{J}}}
\def\sK{{\mathbb{K}}}
\def\sL{{\mathbb{L}}}
\def\sM{{\mathbb{M}}}
\def\sN{{\mathbb{N}}}
\def\sO{{\mathbb{O}}}
\def\sP{{\mathbb{P}}}
\def\sQ{{\mathbb{Q}}}
\def\sR{{\mathbb{R}}}
\def\sS{{\mathbb{S}}}
\def\sT{{\mathbb{T}}}
\def\sU{{\mathbb{U}}}
\def\sV{{\mathbb{V}}}
\def\sW{{\mathbb{W}}}
\def\sX{{\mathbb{X}}}
\def\sY{{\mathbb{Y}}}
\def\sZ{{\mathbb{Z}}}

% Entries of a matrix
\def\emLambda{{\Lambda}}
\def\emA{{A}}
\def\emB{{B}}
\def\emC{{C}}
\def\emD{{D}}
\def\emE{{E}}
\def\emF{{F}}
\def\emG{{G}}
\def\emH{{H}}
\def\emI{{I}}
\def\emJ{{J}}
\def\emK{{K}}
\def\emL{{L}}
\def\emM{{M}}
\def\emN{{N}}
\def\emO{{O}}
\def\emP{{P}}
\def\emQ{{Q}}
\def\emR{{R}}
\def\emS{{S}}
\def\emT{{T}}
\def\emU{{U}}
\def\emV{{V}}
\def\emW{{W}}
\def\emX{{X}}
\def\emY{{Y}}
\def\emZ{{Z}}
\def\emSigma{{\Sigma}}

% entries of a tensor
% Same font as tensor, without \bm wrapper
\newcommand{\etens}[1]{\mathsfit{#1}}
\def\etLambda{{\etens{\Lambda}}}
\def\etA{{\etens{A}}}
\def\etB{{\etens{B}}}
\def\etC{{\etens{C}}}
\def\etD{{\etens{D}}}
\def\etE{{\etens{E}}}
\def\etF{{\etens{F}}}
\def\etG{{\etens{G}}}
\def\etH{{\etens{H}}}
\def\etI{{\etens{I}}}
\def\etJ{{\etens{J}}}
\def\etK{{\etens{K}}}
\def\etL{{\etens{L}}}
\def\etM{{\etens{M}}}
\def\etN{{\etens{N}}}
\def\etO{{\etens{O}}}
\def\etP{{\etens{P}}}
\def\etQ{{\etens{Q}}}
\def\etR{{\etens{R}}}
\def\etS{{\etens{S}}}
\def\etT{{\etens{T}}}
\def\etU{{\etens{U}}}
\def\etV{{\etens{V}}}
\def\etW{{\etens{W}}}
\def\etX{{\etens{X}}}
\def\etY{{\etens{Y}}}
\def\etZ{{\etens{Z}}}

% The true underlying data generating distribution
\newcommand{\pdata}{p_{\rm{data}}}
% The empirical distribution defined by the training set
\newcommand{\ptrain}{\hat{p}_{\rm{data}}}
\newcommand{\Ptrain}{\hat{P}_{\rm{data}}}
% The model distribution
\newcommand{\pmodel}{p_{\rm{model}}}
\newcommand{\Pmodel}{P_{\rm{model}}}
\newcommand{\ptildemodel}{\tilde{p}_{\rm{model}}}
% Stochastic autoencoder distributions
\newcommand{\pencode}{p_{\rm{encoder}}}
\newcommand{\pdecode}{p_{\rm{decoder}}}
\newcommand{\precons}{p_{\rm{reconstruct}}}

\newcommand{\laplace}{\mathrm{Laplace}} % Laplace distribution

\newcommand{\E}{\mathbb{E}}
\newcommand{\Ls}{\mathcal{L}}
\newcommand{\R}{\mathbb{R}}
\newcommand{\emp}{\tilde{p}}
\newcommand{\lr}{\alpha}
\newcommand{\reg}{\lambda}
\newcommand{\rect}{\mathrm{rectifier}}
\newcommand{\softmax}{\mathrm{softmax}}
\newcommand{\sigmoid}{\sigma}
\newcommand{\softplus}{\zeta}
\newcommand{\KL}{D_{\mathrm{KL}}}
\newcommand{\Var}{\mathrm{Var}}
\newcommand{\standarderror}{\mathrm{SE}}
\newcommand{\Cov}{\mathrm{Cov}}
% Wolfram Mathworld says $L^2$ is for function spaces and $\ell^2$ is for vectors
% But then they seem to use $L^2$ for vectors throughout the site, and so does
% wikipedia.
\newcommand{\normlzero}{L^0}
\newcommand{\normlone}{L^1}
\newcommand{\normltwo}{L^2}
\newcommand{\normlp}{L^p}
\newcommand{\normmax}{L^\infty}

\newcommand{\parents}{Pa} % See usage in notation.tex. Chosen to match Daphne's book.

\DeclareMathOperator*{\argmax}{arg\,max}
\DeclareMathOperator*{\argmin}{arg\,min}

\DeclareMathOperator{\sign}{sign}
\DeclareMathOperator{\Tr}{Tr}
\let\ab\allowbreak


%%% new
\newcommand{\diag}{\mathop{\mathrm{diag}}\nolimits}
\usepackage{cancel}
\newcommand{\ttt}[1]{\texttt{#1}}

\newif\ifexercise
\exercisetrue
%\exercisefalse

\newif\ifsolution
\solutiontrue
%\solutionfalse

\usepackage{booktabs}
\usepackage[ruled,vlined]{algorithm2e}

\usepackage{amsthm}
\theoremstyle{definition}
\newtheorem{exercise}{Question}%[chapter]
\newtheorem{answer}{Answer} % asterisk to remove ordering

\newcommand{\Exercise}[1]{
\ifexercise#1\fi
}

\definecolor{answer}{rgb}{0.00, 0.12, 0.60}
\newcommand{\Answer}[1]{
\ifsolution\color{answer}\begin{answer}#1\end{answer}\fi
}

\newif\ifexercise
\exercisetrue
%\exercisefalse

\newif\ifsolution
\solutiontrue
%\solutionfalse

\usepackage{enumitem}
\newcommand{\staritem}{
\addtocounter{enumi}{1}
\item[$\phantom{x}^{*}$\theenumi]}
\setlist[enumerate,1]{leftmargin=*, label=\arabic*.}

\newcommand{\pd}[2]{\dfrac{\partial #1}{\partial #2}}
\newcommand{\relu}{\mathrm{ReLU}}
\newcommand{\vect}[1]{\bm{#1}}

\begin{document}


\fancyhead{}
\fancyfoot{}

\fancyhead[L]{
  \begin{tabular}[b]{l}
    IFT6135-A2026  \\
    Prof: Aishwarya Agrawal \\
  \end{tabular}
}
\fancyhead[R]{
  \begin{tabular}[b]{r}
    Assignment 1  \\
    MLPs, CNNs, Activation Functions, Gradient Descent  \\
  \end{tabular}
}
\fancyfoot[C]{- Do not distribute -}

\vspace{1cm}

\shorthandoff{:}
{\textbf{Due Date: October 13th, 2026 at 11:00 pm}}\\


\vspace{-0.5cm}
\underline{Instructions}%
\renewcommand{\labelitemi}{\textbullet}

\begin{itemize}
% \item \emph{This assignment serves as a warm-up for the following assignments. You are not obliged to finish this assignment, but some of the results here might be useful for the upcoming assignments. Unless otherwise specified, you may use the results in this assignment directly in your answer in the future.}
\item \emph{Use LaTeX and the template we provide when writing your answers.
You may reuse most of the notation shorthands, equations and/or tables.
See the assignment policy on the course website for more details.}
% \item \emph{You will be using Gradescope, you should have received an invitation email (sent to the email address you use for StudiuM). Otherwise reply to the thread named "Gradescope Sign Up" on Piazza to let the TA know your “name”, “email” and “student ID” (matricule) as shown on StudiuM to ask the TA to add you on Gradescope. 
% }
\item \emph{Submit this assignment on Gradescope. Make sure that you select the pages containing the answer for each question on Gradescope. We will not grade questions without selected pages.}
\item \emph{TAs for this assignment are \textbf{Aneri Muni and Serena Chan}.}
\end{itemize}
\vspace{0.2cm}

% Q1
\begin{exercise}[2-2-5-2-2-2]
\Exercise{\label{ex:forward-backward-pass}
Assume you want to train a two-layer MLP with cross-entropy loss for binary classification using gradient descent.

\noindent \textbf{Network architecture:}
\begin{itemize}[nosep]
    \item Input layer: 2 neurons
    \item Hidden layer: 3 neurons with ReLU activation
    \item Output layer: 1 neuron with Sigmoid ($\sigma$) activation
\end{itemize}

\noindent \textbf{The forward pass equations:}
\begin{align}
    \vect{z}^{(1)} &= W_1 \vect{x} + \vect{b}_1 \label{eq:z1} \\
    \vect{h}^{(1)} &= \relu(\vect{z}^{(1)}) \label{eq:h1} \\
    z^{(2)} &= W_2^\top \vect{h}^{(1)} + b_2 \label{eq:z2} \\
    \hat{y} &= \sigmoid(z^{(2)}) = \frac{1}{1 + e^{-z^{(2)}}} \label{eq:yhat}
\end{align}

\noindent The following \textbf{component-wise notations} can be used to expand \eqref{eq:z2}:
\begin{equation*}
    z^{(2)} = W_{2,1}\,h_1^{(1)} + W_{2,2}\,h_2^{(1)} + W_{2,3}\,h_3^{(1)} + b_2
    \label{eq:z2_expanded}
\end{equation*}

\noindent \textbf{Initial parameter values:}
$ \quad
W_1 = \begin{bmatrix} 1 & 0 \\ 1 & 1 \\ 0 & 1 \end{bmatrix}, \quad
\vect{b}_1 = \begin{bmatrix} 0 \\ 0 \\ 0 \end{bmatrix}, \quad
W_2 = \begin{bmatrix} 1 \\ -1 \\ 1 \end{bmatrix}, \quad
b_2 = 0.
$

\noindent \textbf{Training sample:}
$ \quad \biggl\{
\vect{x} = \begin{bmatrix} 1 \\ 1 \end{bmatrix},\; y = 1
\biggr\},
$
where $x$ is the input and $y$ is the label. 

\noindent \textbf{Conventions for this question:}
\begin{itemize}[nosep]
    \item Approximate $e \approx 3$.
    \item Assume the convention $\relu'(0) = 0$.
\end{itemize}

\textbf{Forward pass:} \\
\noindent (1) We will first do a forward pass using the given parameter initialization to compute the network label predictions $\hat{y}$ for each data point, prior to any training. Fill in the values of the given table :
\begin{center}
\renewcommand{\arraystretch}{1.4}
\begin{tabular}{l|c}
    \toprule
    \textbf{Quantity} & $\vect{x} = [1,1]^\top,\; y = 1$ \\
    \midrule
    $\vect{z}^{(1)}$ &   \\
    $\vect{h}^{(1)}$ &   \\
    $z^{(2)}$        &   \\
    $\hat{y}$        &   \\
    \bottomrule
\end{tabular}
\end{center}

\textbf{Backward pass:} \\
The binary cross-entropy loss is given as:
\begin{equation*}
    L = -\Big[ y \ln \hat{y} + (1 - y) \ln(1 - \hat{y}) \Big].
    \label{eq:bce}
\end{equation*}
What is the derivative of the cross-entropy loss function w.r.t $z^{(2)}$, i.e. give the value of $\pd{L}{z^{(2)}}$.

\noindent (3) Using chain rule, provide the analytical formula for:
\begin{equation*}
    \pd{L}{b_2}, \pd{L}{W_2}, \pd{L}{h^{(1)}}, \pd{L}{b_1}, \pd{L}{W_1}. 
\end{equation*}

\noindent (4) Plug in the values for all the parameters and data points to compute the numerical values for:
\begin{equation*}
    \pd{L}{b_2}, \pd{L}{W_2}, \pd{L}{b_1}, \pd{L}{W_1}. 
\end{equation*}

\textbf{Gradient Descent:}\\
\noindent (5) We now apply the vanilla gradient descent rule with learning rate $\eta = 0.5$. Using the update rule for each parameter $\theta$:
\begin{equation}
    \theta^{\text{new}} = \theta^{\text{old}} - \eta\ \cdot \pd{L}{\theta},
    \label{eq:update}
\end{equation}
compute the numerical values of $b^{\text{new}}_1, W^{\text{new}}_1, b^{\text{new}}_2, W^{\text{new}}_2$ after one gradient step.

\noindent (6) Finally, recompute the prediction label for the given data point under the updated parameters. Did the gradient update move the prediction in the right direction?
}

\Answer{
% Write your answer here
}

\end{exercise}

% Q2
\begin{exercise}[2-5-3] (\textbf{Activation Functions})
\Exercise{\label{ex:activation-fxns}

Consider an $L$-layer MLP trained for a multi-class classification task. Let $x \in \mathbb{R}^d$ denote the input vector. For each layer $l \in \{1, \dots, L\}$, let $W^{(l)}$ and $b^{(l)}$ denote the weight matrix and bias vector, respectively. The final layer outputs logits, to which the softmax function is applied, produces the predicted probability distribution, $\hat{y}$. The network is trained to minimize the cross-entropy loss $\mathcal{L}$.

\begin{enumerate}
    \item Assume the activation function is the identity function, $g(a) = a$, for all hidden layers.
    \begin{enumerate}
        \item Recall that a layer with a linear activation function strictly performs an affine transformation. Given two consecutive layers $l$ and $l+1$, prove that their composition collapses into a single equivalent transformation.
        \item Using your answer from (a), justify why an $L$-layer network with linear activations provides no more representational capacity than a single-layer linear classifier.
    \end{enumerate}

    \item Assume all hidden layers use the sigmoid activation function, $g(a) = \sigma(a) = \frac{1}{1 + e^{-a}}$.
    \begin{enumerate}
        \item Determine the absolute maximum value of $g'(a)$.
        \item Using the chain rule, prove an upper bound for the absolute gradient at the first layer, $\left\vert{}\frac{\partial L}{\partial a^{(1)}}\right\vert{}$, in terms of $W_{\max}$, $L$, and the loss gradient at the output layer. Assume that for this question, the hidden layers have a width of 1 (1 neuron per layer, thus weights $w^{(l)}$ and pre-activations $a^{(l)}$ are scalars), and that the norm of every weight matrix is bounded such that $ |w^{(l)} | \leq W_{\max}$ for all $l$.
        \item Using the bound derived in (b), express the mathematical condition on $W_{\max}$ that causes the gradients to vanish exponentially as network depth $L$ increases.
    \end{enumerate}

    \item Assume all hidden layers use the ReLU activation function, $g(a) = \max(0, a)$.
    \begin{enumerate}
        \item Write the expression for the derivative of the ReLU function, $g'(a)$.
        \item If a neuron $i$ in layer $l$ has a pre-activation $a_i^{(l)} \leq 0$, what is its corresponding weight gradient vector?
        \item Compare and contrast how LeakyReLU and Maxout networks can be used to modify the activation step to prevent this.
    \end{enumerate}
\end{enumerate}

}

\Answer{
% Write your answer here
}

\end{exercise}

% Q3
\begin{exercise}[10] (\textbf{Gradient Descent)}
\Exercise{\label{ex:loss-hessian}

Suppose the loss function takes the following form
$$
L=L_0+\frac{1}{2}\left(\mathbf{w}-\mathbf{w}^{\star}\right)^{\mathrm{T}} \mathbf{H}\left(\mathbf{w}-\mathbf{w}^{\star}\right)
$$

where $\mathbf{w}^{\star}$ is the minimum, and the Hessian matrix $\mathbf{H}$ is positive definite and constant. Suppose the initial weight vector $\mathbf{w}^{(0)}$ is chosen to be at the origin and is updated using simple gradient descent. 
$$
\mathbf{w}^{(k)}=\mathbf{w}^{(k-1)}-\rho \nabla L
$$

where $k$ denotes the step number, and $\rho$ is the learning rate (which is assumed to be small). Show using mathematical induction that, after $k$ steps, the components of the weight vector parallel to the eigenvectors of $\mathbf{H}$ can be written as 

$$
w_j^{(k)}=\left\{1-\left(1-\rho \eta_j\right)^k\right\} w_j^{\star}
$$

where $w_j=\mathbf{w}^{\mathrm{T}} \mathbf{u}_j$, and $\mathbf{u}_j$ and $\eta_j$ are the eigenvectors and eigenvalues, respectively, of $\mathbf{H}$ so that

$$
\mathbf{H} \mathbf{u}_j=\eta_j \mathbf{u}_j
$$
 
Show that as $k \rightarrow \infty$, this gives $\mathbf{w}^{(k)} \rightarrow \mathbf{w}^{\star}$ as expected, provided $\left|1-\rho \eta_j\right|<1$.

}

\Answer{
% Write your answer here
}

\end{exercise}


% Q4
\begin{exercise}[15]{
\label{ex:convolutional-neural-networks}
Convolutional Neural Networks
\begin{enumerate}
\item \textbf{(2 points)} Suppose I have a CNN with the following layers in exact order without any additional componenets between them:
\begin{itemize}
    \item A $3\times 3$ convolutional layer with stride $3$.
    \item A $2\times 2$ average pooling layer with stride $2$.
    \item A $2\times 2$ convolutional layer with stride $2$.
\end{itemize}
Give the exact single layer that would produce the same output as this sequence of layers. Write out a formula that describes the weights for the layer.

\item \textbf{(4 points)} Given the CNN defined by the layers below, fill in the shape of the output volume and the number of parameters at each layer in the following table (Table \ref{tab:cnn-table}). Write the shapes in the format $(H, W, C)$, where $H$, $W$, $C$ are the height, width and channel dimensions. Unless specified, assume padding 1 and stride 1 where appropriate. Layers are expressed as follows:
\begin{itemize}
    \item \textsc{Conv(H,C)}: A convolution layer with filters of size $H\times W$ with $C$ channels with appropriate weights and biases.
    \item \textsc{ReLU}: ReLU activation.
    \item \textsc{MaxPooling(N)}: $N\times N$ max pooling.
    \item \textsc{FC(D)}: A fully connected layer with $D$ outputs with appropriate weights and biases.
    \item \textsc{BatchNorm}: Batch normalization.
    \item \textsc{Reshape}: A flattening layer.
\end{itemize}
\begin{table}[h]
    \centering
    \begin{tabular}{l|c|c}
    \toprule
    Layer & Output Dimensions & Number of Parameters \\
    \midrule
    \textsc{Input} & $64\times 64 \times 3$ & $0$\\
    \textsc{Conv(4,8)}, stride = 2 & & \\
    \textsc{BatchNorm} & & \\
    \textsc{ReLU} & & \\
    \textsc{MaxPooling(2)} & & \\
    \textsc{Conv(3,16)} & & \\
    \textsc{BatchNorm} & & \\
    \textsc{ReLU} & & \\
    \textsc{MaxPooling(4)} & & \\
    \textsc{Reshape} & & \\
    \textsc{FC(10)} & & \\
    \bottomrule
    \end{tabular}
    \caption{Table to fill out for part 2.}
    \label{tab:cnn-table}
\end{table}

\item \textbf{(2 points)} In the above setup, there are specific parameters that can be removed without affecting performance. What are they and why can they be removed?

\item \textbf{(2 points)} Fill out the following table (Table \ref{tab:receptive-field-table}) by calculating the receptive field of a single pixel at every layer output. You can assume each layer has only a single channel and stride is $1$ for each layer.

\begin{table}[h]
    \centering
    \begin{tabular}{l|c}
    \toprule
    Layer & Receptive Field \\
    \midrule
    \textsc{Conv(5,1)} with stride 3 & \\
    \textsc{MaxPooling(2)} with stride 2 & \\
    \textsc{BatchNorm} & \\
    \textsc{Conv(4,1)} with stride 2 & \\
    \textsc{MaxPooling(2)} with stride 1 & \\
    \textsc{Conv(3,1)} with stride 1& \\
    \bottomrule
    \end{tabular}
    \caption{Table to fill out for part 4.}
    \label{tab:receptive-field-table}
\end{table}

\item \textbf{(4 points)} Suppose your input is a black and white image and consists only of pixels with values 0 or 1, where 1 is white and 0 is black (assume the image consists only of a single channel). Design two $3\times 3$ convolution kernels, one that detects left-to-right dark-to-light vertical boundaries and another that detects left-to-right light-to-dark vertical boundaries. Your kernels should be as simple as possible (for simplicity, restrict values to be from -1 to 1), with output 1 for areas around the edges and 0 elsewhere. Explain what construct enables your kernels to detect these boundaries and not other boundaries (such as horizontal ones).
\end{enumerate}
}

\Answer{
% Write your answer here
}

\end{exercise}



\end{document}
